% Example section for the Humber ASV template.
% This section demonstrates all major features of the template.
% It is meant to be included temporarily for reference and removed before final submission.
% MIT License
%
% Copyright (c) 2026 HumberASV
% Copyright (c) 2026 Carson Fujita
%
% Permission is hereby granted, free of charge, to any person obtaining a copy
% of this software and associated documentation files (the "Software"), to deal
% in the Software without restriction, including without limitation the rights
% to use, copy, modify, merge, publish, distribute, sublicense, and/or sell
% copies of the Software, and to permit persons to whom the Software is
% furnished to do so, subject to the following conditions:
%
% The above copyright notice and this permission notice shall be included in all
% copies or substantial portions of the Software.
%
% THE SOFTWARE IS PROVIDED "AS IS", WITHOUT WARRANTY OF ANY KIND, EXPRESS OR
% IMPLIED, INCLUDING BUT NOT LIMITED TO THE WARRANTIES OF MERCHANTABILITY,
% FITNESS FOR A PARTICULAR PURPOSE AND NONINFRINGEMENT. IN NO EVENT SHALL THE
% AUTHORS OR COPYRIGHT HOLDERS BE LIABLE FOR ANY CLAIM, DAMAGES OR OTHER
% LIABILITY, WHETHER IN AN ACTION OF CONTRACT, TORT OR OTHERWISE, ARISING FROM,
% OUT OF OR IN CONNECTION WITH THE SOFTWARE OR THE USE OR OTHER DEALINGS IN THE
% SOFTWARE.

\section{Template Usage Guide}

This section documents every feature of the Humber ASV \LaTeX{} template.
Remove or comment out this file from \texttt{main.tex} before submitting.

\subsection{Project Structure}

The repository is organized as follows:

\begin{longtable}{ll}
    \textbf{Path} & \textbf{Purpose} \\
    \hline
    \texttt{main.tex} & Entry point; orchestrates the document \\
    \texttt{preamble.sty} & All packages, settings, and custom commands \\
    \texttt{sections/} & One \texttt{.tex} file per content section \\
    \texttt{citations/citations.bib} & BibTeX bibliography entries \\
    \texttt{images/} & Figures and the logo (\texttt{logo.png}) \\
    \texttt{code/} & Source code snippets or listings \\
\end{longtable}

\subsection{Title Page}

The \texttt{\textbackslash asvtitle} command generates a centred title page with the logo,
title, institution name, version, and date.

\begin{verbatim}
% Optional date defaults to \today when omitted
\asvtitle{My Document Title}{1.0.0}
\asvtitle[2026-01-15]{My Document Title}{1.0.0}
\end{verbatim}

\subsection{Table of Contents}

Call \texttt{\textbackslash startcontents} immediately after \texttt{\textbackslash asvtitle}.
It generates the table of contents, list of figures, and list of tables on separate pages.

\begin{verbatim}
\startcontents
\end{verbatim}

\subsection{Adding Sections}

Create a new \texttt{.tex} file inside \texttt{sections/} and include it in \texttt{main.tex}
using \texttt{\textbackslash input}:

\begin{verbatim}
\input{sections/introduction.tex}
\input{sections/methodology.tex}
\input{sections/results.tex}
\end{verbatim}

Each call to \texttt{\textbackslash section\{\}} inside a section file automatically appears in
the table of contents.

\subsection{Appendix}

Call \texttt{\textbackslash startappendix} before any appendix \texttt{\textbackslash input} calls.
Section numbers switch to letters (A, B, C, \ldots) automatically.

\begin{verbatim}
\startappendix
% Glossary for the document
% Optional, but you must delete the call in main.tex
% MIT License
%
% Copyright (c) 2026 HumberASV
% Copyright (c) 2026 Carson Fujita
%
% Permission is hereby granted, free of charge, to any person obtaining a copy
% of this software and associated documentation files (the "Software"), to deal
% in the Software without restriction, including without limitation the rights
% to use, copy, modify, merge, publish, distribute, sublicense, and/or sell
% copies of the Software, and to permit persons to whom the Software is
% furnished to do so, subject to the following conditions:
%
% The above copyright notice and this permission notice shall be included in all
% copies or substantial portions of the Software.
%
% THE SOFTWARE IS PROVIDED "AS IS", WITHOUT WARRANTY OF ANY KIND, EXPRESS OR
% IMPLIED, INCLUDING BUT NOT LIMITED TO THE WARRANTIES OF MERCHANTABILITY,
% FITNESS FOR A PARTICULAR PURPOSE AND NONINFRINGEMENT. IN NO EVENT SHALL THE
% AUTHORS OR COPYRIGHT HOLDERS BE LIABLE FOR ANY CLAIM, DAMAGES OR OTHER
% LIABILITY, WHETHER IN AN ACTION OF CONTRACT, TORT OR OTHERWISE, ARISING FROM,
% OUT OF OR IN CONNECTION WITH THE SOFTWARE OR THE USE OR OTHER DEALINGS IN THE
% SOFTWARE.

\AtBeginEnvironment{thebibliography}{\renewcommand{\autocite}{\textcite}} % Change the default citation command to textcite for better readability in the bibliography
% right now the default is set to numeric, but you can change it to authoryear or any other style supported by biblatex. 
% Adjust the citation command accordingly if you change the style. 
% For example, if you switch to authoryear, you might want to use \cite instead of \textcite in the bibliography.
% \AtBeginEnvironment{thebibliography}{\renewcommand{\autocite}{\cite}} % Uncomment this line if you prefer longer citations instead

% Note: Make sure you add the relevant entries to your .bib file for the citations used in the glossary.
\section{Glossary}
\begin{description}
	\item[\textbf{Autonomous Surface Vehicle (ASV)}]	
	An autonomous surface vehicle (ASV) is an unmanned vessel that operates on the sea surface without real-time input or control from human operators.\autocite{TREMBANIS2021251}

	\item[\textbf{Base Station}]	
	A fixed radio communications station that can communicate with mobile radio devices in its neighbourhood. Forms of radio communication that have base stations include radio wireless data and mobile telephony...\autocite{BASESTATION2018}

	\item[\textbf{Loon-E}]	
	Humber ASV's current ASV.\autocite{HuMBERASVTEAM}
	
	\item[\textbf{HumberASV}]	
	Humber Polytechnic's Mechatronics club. A multidisciplinary team of passionate Humber Polytechnic students and professionals dedicated to advancing autonomous maritime technology.\autocite{HUMBERASVTEAM}

	\item[\textbf{LIDAR}]
	LiDAR, an acronym for “light detection and ranging,” is a remote-sensing technology that uses laser beams to measure precise distances and movement in an environment, in real time.\autocite{LIDAR}

	\item[\textbf{Robot Operating System (ROS)}]	
	The Robot Operating System (ROS) is a set of software libraries and tools that help you build robot applications.\autocite{ROS}
\end{description}

\section{Bibliography}
\printbibliography[heading=none]
\end{verbatim}

\subsection{Citations}

Add BibTeX entries to \texttt{citations/citations.bib}, then cite inline with
\texttt{\textbackslash autocite\{key\}}. Example entry and usage:

\begin{verbatim}
% citations/citations.bib
@misc{example2026,
    author = {Author, A.},
    title  = {Example Reference},
    year   = {2026},
    url    = {https://example.com}
}
\end{verbatim}

\begin{verbatim}
% In a section file
The system was validated against prior work \autocite{example2026}.
\end{verbatim}

Run \texttt{pdflatex}, then \texttt{bibtex}, then \texttt{pdflatex} twice to resolve all references.

\subsection{Figures}

Place image files in \texttt{images/}. Use the \texttt{[H]} specifier (provided by the
\texttt{float} package) to pin a figure exactly where it appears in the source.

\begin{verbatim}
\begin{figure}[H]
    \centering
    \includegraphics[width=0.6\linewidth]{./images/my-diagram.png}
    \caption{A descriptive caption.}
    \label{fig:my-diagram}
\end{figure}

See Figure~\ref{fig:my-diagram} for details.
\end{verbatim}

\subsection{Tables}

Use \texttt{longtable} for multi-page tables, \texttt{multirow} for spanning rows, and
\texttt{makecell} for line breaks inside cells.

\begin{verbatim}
\begin{longtable}{|l|c|p{6cm}|}
    \hline
    \textbf{Name} & \textbf{Value} & \textbf{Description} \\
    \hline
    Speed & 2 m/s & Maximum forward velocity \\
    \makecell{Turn \\ radius} & 1.5 m & Minimum turning radius at full speed \\
    \hline
    \caption{System parameters.}
    \label{tab:params}
\end{longtable}
\end{verbatim}

\subsection{Landscape Pages}

Wrap wide content in the \texttt{landscape} environment from the \texttt{pdflscape} package.
Only the enclosed pages rotate; the rest of the document is unaffected.

\begin{verbatim}
\begin{landscape}
    % Wide table or figure here
\end{landscape}
\end{verbatim}

\subsection{Callout Boxes}

Three callout box commands are available for drawing attention to important content.

\subsubsection*{Warning Box}

Use \texttt{\textbackslash warningbox} for critical information the reader must not miss.

\warningbox{This is a warning. Ignoring this may result in hardware damage or unsafe operation.}

\begin{verbatim}
\warningbox{This is a warning. Ignoring this may result in hardware damage.}
\end{verbatim}

\subsubsection*{Note Box}

Use \texttt{\textbackslash notebox} for optional guidance or contextual information.

\notebox{This is a note. It provides additional context that is helpful but not critical.}

\begin{verbatim}
\notebox{This is a note. It provides additional context.}
\end{verbatim}

\subsubsection*{Info Box}

Use \texttt{\textbackslash infobox} for positive status updates or useful reference details.

\infobox{This is an info box. Use it to highlight useful facts or confirmed status.}

\begin{verbatim}
\infobox{This is an info box. Use it to highlight useful facts.}
\end{verbatim}

All three commands accept an optional second argument for supplementary content below the main message:

\begin{verbatim}
\warningbox{Primary warning message.}[See Section~\ref{sec:safety} for details.]
\end{verbatim}
